\documentclass[10pt,a4paper]{amsart}
\usepackage[margin=19mm]{geometry}
\usepackage{amsmath,amssymb,mathtools}
\usepackage[scr=rsfs,cal=euler]{mathalfa}
\usepackage{xfrac}
\usepackage[unicode,hidelinks,bookmarks=false]{hyperref}
\newcommand{\ie}{i.e.\ }
\newcommand{\cf}{cf.\ }
\newcommand{\resp}{respectively\ }
\newcommand{\Subsection}{\subsection}
\providecommand{\nobreakdash}{\nobreak\hbox{-}}
\setlength{\emergencystretch}{2em}
\multlinegap0pt
\usepackage{amssymb}
\usepackage[shortlabels]{enumitem}
\usepackage{tikz}
\usepackage{xcolor}
\newtheorem{prop}{Proposition}
\title{On some representations of Metacyclic groups whose integral forms can be computed from a single residual representation}
\author{Aguil\'o-Vidal, Bruno}
\date{Source: arXiv:2608.08854v1 (CC BY 4.0)}
\begin{document}
\maketitle

\noindent\emph{Source and licence.} On some representations of Metacyclic groups whose integral forms can be computed from a single residual representation. Version arXiv:2608.08854v1 (CC BY 4.0), distributed by arXiv under the Creative Commons Attribution 4.0 International licence, http://creativecommons.org/licenses/by/4.0/, as recorded in the arXiv licence metadata for this exact version and shown on the abstract page https://arxiv.org/abs/2608.08854v1. Full licence notice: \texttt{licenses/arxiv-2608.08854-CC-BY-4.0.txt}.\par\smallskip

\noindent\textit{Editorial scope.} Counted unit: the article's prop p44 (source0.tex lines 792--798). The article's complete original proof of it is delivered with it (source0.tex lines 800--847, one \texttt{proof} environment of 48 lines). No mathematics is written, repaired or re-derived: the statement and the proof are the article's own lines, in the article's own notation, and no retained line is substituted or re-flowed.  No further source-local context is needed: every cross-reference of every retained line resolves inside the two delivered spans.  Nothing was omitted from any retained span, so the package carries no omission record.  No retained line contains a citation, so the wrapper carries no bibliography and no bibliography entry.  No retained line contains a figure environment, an image-inclusion command or drawing code, so no image is delivered and the package declares no asset dependency.  Editorial formatting only, in addition: an AMS \texttt{amsart} preamble that restates the article's own declarations and macros used by the retained lines, namely the environments \texttt{prop} on separate counters, together with the standard packages those lines need.  The retained environments print in this document as Proposition 1 in order of appearance, whereas the article prints the same items with the numbering of its own full document.  The complete native source of the article as distributed in the arXiv e-print for this exact version is bundled unchanged as \texttt{sources/207198/source0.tex}.
\begin{prop}\label{p44}
    Let $\phi:G\rightarrow\mathrm{GL}_n(K)$ be an absolutely irreducible 
    representation of a finite group $G$. Then there is an $h_k(n)$-to-one
    map from the set of integral forms of $\phi$ to the set of
    $\mathcal{O}_K$-maximal orders containing $\phi(G)$ that are 
    $\mathrm{GL}_n(K)$-conjugate to $\mathbb{M}_n(\mathcal{O}_K)$.
\end{prop}

\begin{proof}
    Let $\phi'$ be a conjugate of $\phi$, say 
    $\phi'(g)=\mathbf{T}\phi(g)\mathbf{T}^{-1}$ , where 
    $\mathbf{T}\in\mathrm{GL}_n(K)$. Then $\phi'(G)\subseteq
    \mathrm{GL}_n(\mathcal{O}_K)$, or equivalently, $\phi'(G)\subseteq
    \mathbb{M}_n(\mathcal{O}_K)$, precisely when $\phi(G)\subseteq
    \mathbf{T}^{-1}\mathbb{M}_n(\mathcal{O}_K)\mathbf{T}$. Furthermore,
    $\mathbf{T}^{-1}\mathbb{M}_n(\mathcal{O}_K)\mathbf{T}=
    \mathbf{S}^{-1}\mathbb{M}_n(\mathcal{O}_K)\mathbf{S}$, precisely when
    $\mathbf{N}=\mathbf{S}\mathbf{T}^{-1}$ belongs to the normalizer $N$ 
    of the order $\mathbb{M}_n(\mathcal{O}_K)$. In other words, 
    two integral representations 
    $\phi'(g)=\mathbf{T}\phi(g)\mathbf{T}^{-1}$ and
    $\phi''(g)=\mathbf{S}\phi(g)\mathbf{S}^{-1}$ correspond to the same
    maximal order if and only if they satisfy 
    $\phi''(g)=\mathbf{N}\phi'(g)\mathbf{N}^{-1}$ for an element 
    $\mathbf{N}\in N$. This means that $N$ act on the set of 
    representations defining the same maximal order by conjugation.
    Since the normalizer of any absolutely irreducible representation 
    is the group $K^*$ of scalar matrices, the representations $\phi'$
    and $\phi''$ are $\mathrm{GL}_n(\mathcal{O}_K)$-conjugates if and
    only if $\mathbf{N}\in K^*\mathrm{GL}_n(\mathcal{O}_K)$. 
    It suffices, therefore, to prove that the quotient 
    $N/K^*\mathrm{GL}_n(\mathcal{O}_K)$ has $h_K(n)$ elements.
    
    Note that $\mathbf{N}\in N$ if and only if
    $\mathbf{N} \mathcal{O}_K^n=J\mathcal{O}_K^n$ for some fractional ideal
    $J=J(\mathbf{N})\subseteq K$. This defines a map from $N$ to the ideal 
    group of $K$ whose kernel is the stabilizer of the lattice 
    $\mathcal{O}_K^n$, i.e., the group $\mathrm{GL}_n(\mathcal{O}_K)$.
    We claim that the image $\mathcal{J}$ coincides with the set $\mathcal{J}'$
    of the ideals $J$ whose $n$-th power $J^n$ is principal. This implies that
    $N/K^*\mathrm{GL}_n(\mathcal{O}_K)$ is isomorphic to 
    $\mathcal{G}_K(n)$, and the result follows.

    To prove the claim, we take the determinant of both sides of the identity 
    $\mathbf{N} \mathcal{O}_K^n=J\mathcal{O}_K^n$. We 
    obtain that $\det(\mathbf{N})\in K^*$ generates the ideal $J^n$.
    This proves $\mathcal{J}\subseteq\mathcal{J}'$. To prove 
    $\mathcal{J}'\subseteq\mathcal{J}$, we observe that, if $J^n$ is principal, 
    the lattice $J\mathcal{O}_K^n$ is free since its determinant class is 
    principal, and therefore there is a matrix satisfying 
    $\mathbf{N} \mathcal{O}_K^n=J\mathcal{O}_K^n$. Furthermore,
    $\mathbf{N}\mathbb{M}_n(\mathcal{O}_K)\mathbf{N}^{-1}$ is the endomorphism
    ring of the latter lattice, and therefore it coincides with
    $\mathbb{M}_n(\mathcal{O}_K)$. We conclude that $\mathbf{N}\in N$.
    The result follows.
\end{proof}

\end{document}
